\documentclass[11pt,a4paper]{article}

\usepackage[margin=1in]{geometry}
\usepackage{amsmath,amssymb,amsthm}
\usepackage{mathtools}
\usepackage{hyperref}
\usepackage{cleveref}
\usepackage[ruled,vlined,linesnumbered]{algorithm2e}

\newtheorem{theorem}{Theorem}[section]
\newtheorem{lemma}[theorem]{Lemma}
\newtheorem{proposition}[theorem]{Proposition}
\newtheorem{corollary}[theorem]{Corollary}
\newtheorem{definition}[theorem]{Definition}
\newtheorem{remark}[theorem]{Remark}

\DeclareMathOperator{\proj}{proj}
\DeclareMathOperator{\AM}{AM}
\DeclareMathOperator{\HM}{HM}
\DeclareMathOperator{\BR}{BR}
\DeclareMathOperator{\supp}{supp}
\newcommand{\R}{\mathbb{R}}
\newcommand{\E}{\mathbb{E}}
\newcommand{\Prob}{\mathbb{P}}
\newcommand{\ip}[2]{\langle #1, #2 \rangle}
\newcommand{\norm}[1]{\|#1\|}

\title{Fixed-Point Nash Equilibrium Search with Bayesian Stopping:\\
A Functional-Analytic Approach to Explore-Exploit in Multi-Agent Systems}
\author{Eugene Shcherbinin\\
\textit{London School of Economics}\\
\texttt{y.shcherbinin@lse.ac.uk}}
\date{March 2026}

\begin{document}
\maketitle

\begin{abstract}
We develop a fixed-point theoretic approach to Nash equilibrium (NE)
discovery in multi-agent systems. A Nash equilibrium is a fixed point of
the joint best-response correspondence $\BR: \Delta \rightrightarrows \Delta$,
where $\Delta = \prod_i \Delta(\mathcal{A}_i)$ is the product of strategy
simplices. By smoothing $\BR$ with a temperature parameter $\tau > 0$,
we obtain a continuous map $\BR_\tau: \Delta \to \Delta$ that is a
contraction for sufficiently small $\tau$ (by Banach's fixed-point theorem).
We then pose the problem of finding the \emph{best} NE as a Bayesian
optimization problem over the fixed-point set, with a Good-Turing-style
estimator governing when to stop exploring. Key insight: \emph{all selfish
agents benefit from thorough NE search}, creating emergent cooperation
in the exploration phase even in competitive games. This ``Fixed-Point
NE Search'' (FP-NE) serves as a meta-algorithm wrapping the $\Omega$-gradient,
initializing it at the best discovered equilibrium.
\end{abstract}

\tableofcontents


%% ========================================================================
\section{Setup}
\label{sec:setup}
%% ========================================================================

Consider an $N$-player normal-form game $\mathcal{G} = (N, \{\mathcal{A}_i\}_{i=1}^N, \{R_i\}_{i=1}^N)$
where agent $i$ has action set $\mathcal{A}_i$ with $|\mathcal{A}_i| = d_i$
and payoff function $R_i : \mathcal{A} \to \R$, where $\mathcal{A} = \prod_i \mathcal{A}_i$.

Mixed strategies live on the product simplex
$\Delta = \prod_{i=1}^N \Delta(\mathcal{A}_i)$, where
$\Delta(\mathcal{A}_i) = \{\pi_i \in \R^{d_i}_{\geq 0} : \sum_{a} \pi_i(a) = 1\}$.

The expected payoff to agent $i$ under joint mixed strategy $\pi = (\pi_1, \ldots, \pi_N)$ is:
\begin{equation}
    V_i(\pi) = \sum_{a \in \mathcal{A}} \left(\prod_{j=1}^N \pi_j(a_j)\right) R_i(a).
\end{equation}

\begin{definition}[Best Response]
The \emph{best-response correspondence} for agent $i$ is:
\begin{equation}
    \BR_i(\pi_{-i}) = \arg\max_{\pi_i \in \Delta(\mathcal{A}_i)} V_i(\pi_i, \pi_{-i}).
\end{equation}
The \emph{joint best-response} is $\BR(\pi) = (\BR_1(\pi_{-1}), \ldots, \BR_N(\pi_{-N}))$.
\end{definition}

\begin{definition}[Nash Equilibrium as Fixed Point]
A strategy profile $\pi^* \in \Delta$ is a Nash equilibrium if and only if
$\pi^* \in \BR(\pi^*)$, i.e., $\pi^*$ is a \emph{fixed point} of $\BR$.
\end{definition}


%% ========================================================================
\section{Smoothed Best Response and Contraction}
\label{sec:contraction}
%% ========================================================================

The best-response correspondence $\BR$ is set-valued and discontinuous,
making fixed-point iteration problematic. We replace it with the
\emph{softmax (logit) best response}:

\begin{definition}[Smoothed Best Response]
For temperature $\tau > 0$, define the smoothed best-response map
$\BR_\tau : \Delta \to \Delta$ by:
\begin{equation}\label{eq:soft-br}
    \BR_{\tau,i}(\pi_{-i})(a) = \frac{\exp\left(\frac{1}{\tau} \sum_{a_{-i}} \pi_{-i}(a_{-i}) R_i(a, a_{-i})\right)}
    {\sum_{a'} \exp\left(\frac{1}{\tau} \sum_{a_{-i}} \pi_{-i}(a_{-i}) R_i(a', a_{-i})\right)}.
\end{equation}
\end{definition}

The key property: $\BR_\tau$ is a smooth self-map on $\Delta$, and its
fixed points are the \emph{quantal response equilibria} (QRE) of
McKelvey and Palfrey (1995). As $\tau \to 0$, QRE converge to Nash equilibria.

\begin{theorem}[Contraction Property]\label{thm:contraction}
For any finite game with payoffs bounded by $\|R_i\|_\infty \leq M$,
the smoothed best-response map $\BR_\tau$ satisfies:
\begin{equation}
    \|\BR_\tau(\pi) - \BR_\tau(\pi')\|_\infty \leq \frac{M}{\tau} \cdot \|\pi - \pi'\|_\infty
\end{equation}
for all $\pi, \pi' \in \Delta$. In particular, if $\tau > M$, then
$\BR_\tau$ is a \emph{contraction mapping} with Lipschitz constant $c = M/\tau < 1$.
\end{theorem}

\begin{proof}
The softmax function $\sigma_\tau(x)(a) = \exp(x_a/\tau) / \sum_{a'} \exp(x_{a'}/\tau)$
has Lipschitz constant $1/(2\tau)$ in $\|\cdot\|_\infty$ norm.
The linear map $\pi_{-i} \mapsto R_i(\cdot, \pi_{-i})$ has operator norm bounded by $M$.
By composition, $\BR_{\tau,i}$ is Lipschitz with constant $M/\tau$.
\end{proof}

\begin{corollary}[Banach Fixed-Point Theorem Applied]
When $\tau > M$, the damped iteration
$\pi_{t+1} = (1-\alpha)\pi_t + \alpha \cdot \BR_\tau(\pi_t)$ converges
to the unique QRE at geometric rate $\max(1-\alpha, c)^t$ from any initialization.
\end{corollary}

\begin{remark}
For $\tau \leq M$, $\BR_\tau$ may have multiple fixed points (corresponding
to multiple QRE/NE). This is precisely the regime where the
\emph{equilibrium selection problem} arises and where our Bayesian search
becomes valuable.
\end{remark}


%% ========================================================================
\section{Bayesian NE Discovery}
\label{sec:bayesian}
%% ========================================================================

When $\tau$ is small (close to NE), the game may have multiple fixed
points. We frame the problem of finding the \emph{best} fixed point
as a species-discovery problem.

\subsection{The Explore-Exploit Tradeoff}

\begin{definition}[NE Search Problem]
Given a game $\mathcal{G}$ with unknown set of Nash equilibria
$\mathcal{E} = \{\pi^{*,1}, \ldots, \pi^{*,K}\}$ (where $K$ is unknown),
find $\pi^{*,k}$ maximizing a selection criterion $\phi(\pi^{*,k})$
(e.g., social welfare $\sum_i V_i(\pi^{*,k})$, or agent $i$'s payoff $V_i(\pi^{*,k})$).
\end{definition}

Each ``search'' consists of:
\begin{enumerate}
    \item Sample a random initialization $\pi_0 \sim \text{Dir}(\mathbf{1})$
    (or perturb a known NE).
    \item Run damped $\BR_\tau$ iteration until convergence.
    \item Record the fixed point (if residual $\|\BR_\tau(\pi) - \pi\| < \epsilon$).
    \item Check if it is distinct from previously discovered NE.
\end{enumerate}

\subsection{Good-Turing Estimator for Total NE Count}

We adapt the classical species-discovery framework to NE counting.
After $n$ searches discovering $d$ distinct NE, the probability of
finding a \emph{new} NE on the next search is estimated by:

\begin{equation}\label{eq:coverage}
    \hat{p}_{\text{new}} = 1 - \hat{C}_n, \qquad
    \hat{C}_n = 1 - \frac{f_1}{n}
\end{equation}
where $f_1$ is the number of NE discovered exactly once (singletons)
and $\hat{C}_n$ is the Good-Turing coverage estimator.

The Chao1 estimator for total species (NE) count:
\begin{equation}
    \hat{K} = d + \frac{f_1^2}{2 f_2}
\end{equation}
where $f_2$ is the number of NE discovered exactly twice.

\subsection{Bayesian Stopping Rule}

\begin{definition}[Stopping Criterion]
Define the \emph{exploration value} for agent $i$ as:
\begin{equation}
    \text{EV}_i = \Prob(\exists \text{ undiscovered } \pi^* : V_i(\pi^*) > V_i^{\text{best}})
\end{equation}
where $V_i^{\text{best}} = \max_{k \leq d} V_i(\pi^{*,k})$ is the best
discovered payoff for agent $i$.

Stop exploring when $\text{EV}_i < \delta$ for all agents $i$.
\end{definition}

The probability decomposes as:
\begin{equation}
    \text{EV}_i \leq \Prob(\text{undiscovered NE exist}) \times
    \Prob(\text{random NE is better} \mid \text{it exists}).
\end{equation}

The first factor is estimated by $1 - d/\hat{K}$. The second factor
can be bounded using the empirical distribution of discovered NE payoffs.


%% ========================================================================
\section{Cooperation from Selfish Search}
\label{sec:cooperation}
%% ========================================================================

The central insight: \emph{in games with multiple NE, the search for the
best NE is a public good that selfish agents voluntarily provide.}

\begin{proposition}[Search as Public Good]\label{prop:public-good}
Consider a game with NE set $\mathcal{E} = \{\pi^{*,1}, \ldots, \pi^{*,K}\}$
where $K \geq 2$. Let $\pi^{*,P}$ denote the Pareto-best NE
(maximizing $\sum_i V_i$). If agents use a selection criterion
$\phi_i(\pi) = V_i(\pi)$, then under mild conditions on the payoff
structure\footnote{Specifically, when the correlation between agents'
payoffs across NE is positive: $\text{Corr}(V_i(\pi^{*,k}), V_j(\pi^{*,k})) > 0$ over $k$.},
increasing the search budget benefits \emph{all} agents:
\begin{equation}
    \frac{\partial}{\partial n} \E\left[\max_{k \leq d(n)} V_i(\pi^{*,k})\right] > 0
    \quad \forall i.
\end{equation}
\end{proposition}

\begin{proof}[Sketch]
More search discovers more NE. Across NE, payoffs tend to be positively
correlated (Pareto-better NE are better for most agents). Thus the
probability that a newly discovered NE improves agent $i$'s best
available payoff is positive for all $i$ simultaneously.
\end{proof}

\begin{remark}
This is a \emph{functional-analytic} grounding of cooperation: it
arises not from altruism or mechanism design but from the topological
structure of the fixed-point set. The fixed points of $\BR$ are
shared landmarks that all agents want to map.
\end{remark}


%% ========================================================================
\section{Integration with the $\Omega$-Framework}
\label{sec:omega}
%% ========================================================================

FP-NE search serves as a \emph{Phase 0} preceding the $\Omega$-gradient:

\begin{algorithm}[H]
\SetAlgoLined
\KwIn{Game $\mathcal{G}$, $\Omega$-config $(\gamma, w, \lambda, \beta)$,
FP-NE budget $B$, confidence $\delta$}
\KwOut{Learned policy profile $\hat{\pi}$}

\tcp{Phase 0: Fixed-Point NE Search}
Initialize $\mathcal{D} \leftarrow \emptyset$ \tcp*{discovered NE}
\For{$s = 1, \ldots, B$}{
    Sample $\pi_0 \sim \text{Dir}(\mathbf{1})$ or perturb best known NE\;
    $\pi^* \leftarrow \text{DampedBR}_\tau(\pi_0)$ \tcp*{Banach iteration}
    \If{$\|\BR_\tau(\pi^*) - \pi^*\| < \epsilon$ \textbf{and} $\pi^*$ is new}{
        $\mathcal{D} \leftarrow \mathcal{D} \cup \{\pi^*\}$\;
    }
    \If{$\text{EV}_i < \delta$ for all $i$}{
        \textbf{break} \tcp*{Bayesian stopping}
    }
}
$\hat{\pi}_{\text{init}} \leftarrow \arg\max_{\pi^* \in \mathcal{D}} \sum_i V_i(\pi^*)$\;

\tcp{Phase 1: $\Omega$-Gradient from best NE}
$\pi_0 \leftarrow \hat{\pi}_{\text{init}}$\;
\For{$n = 1, 2, \ldots$}{
    $\pi_{n+1} \leftarrow \proj_\Delta\!\left(\pi_n + \gamma_n \cdot w_n \cdot (\hat{v}_n + \lambda_n \cdot \text{OS}_n + \beta_n \cdot \text{coop}_n)\right)$\;
}
\Return{$\pi_n$}
\caption{$\Omega$-PG with FP-NE Initialization}
\label{alg:omega-fpne}
\end{algorithm}

\begin{theorem}[FP-NE Preserves $\Omega$-Convergence]
If the $\Omega$-gradient converges to a stable NE $\pi^*$ from any
initialization in basin $\mathcal{B}(\pi^*)$, and FP-NE search
discovers $\pi^*$ with probability $\geq 1 - \epsilon$, then
$\Omega$-PG with FP-NE initialization converges to the best
discoverable NE with probability $\geq 1 - \epsilon$.
\end{theorem}

The advantage over random initialization: random init lands in the
basin of an arbitrary NE. FP-NE search explicitly maps the fixed-point
landscape and selects the best basin before starting the $\Omega$-gradient.


%% ========================================================================
\section{Connection to Kakutani and Brouwer}
\label{sec:kakutani}
%% ========================================================================

The existence of Nash equilibria is guaranteed by Kakutani's fixed-point
theorem (for the set-valued BR correspondence) or equivalently by
Brouwer's theorem (for the smoothed map $\BR_\tau$).

\begin{theorem}[Nash 1950, via Kakutani]
Every finite game has at least one mixed-strategy Nash equilibrium.
\end{theorem}

Our contribution is algorithmic: given that fixed points \emph{exist},
we provide an efficient method to \emph{find them all} (or at least
the best ones) with Bayesian confidence guarantees on completeness.

The functional-analytic perspective:
\begin{itemize}
    \item $\Delta$ is compact and convex (finite-dimensional simplex).
    \item $\BR_\tau : \Delta \to \Delta$ is continuous for $\tau > 0$.
    \item By Brouwer, $\BR_\tau$ has at least one fixed point.
    \item For $\tau > M$: unique fixed point (Banach contraction).
    \item For $\tau \leq M$: possibly multiple fixed points.
    Decreasing $\tau$ from $> M$ to near $0$ traces the \emph{QRE
    correspondence} connecting the unique high-$\tau$ QRE to the
    (possibly multiple) NE.
\end{itemize}

This suggests a \emph{homotopy method}: start at high $\tau$ (unique
contraction), gradually decrease $\tau$, and track how the unique
fixed point bifurcates into multiple NE. Branch points in this
homotopy correspond to equilibrium selection events.


%% ========================================================================
\section{Experimental Predictions}
\label{sec:predictions}
%% ========================================================================

\begin{enumerate}
    \item \textbf{NE Discovery}: FP-NE search discovers all NE in games
    with known equilibrium sets (validated against nashpy).
    \item \textbf{Bayesian Stopping}: The stopping criterion correlates
    with actual search completeness.
    \item \textbf{Equilibrium Selection}: FP-NE + $\Omega$-PG achieves
    higher social welfare than $\Omega$-PG alone in games with multiple NE.
    \item \textbf{Contraction Rates}: Empirical contraction constants
    of $\BR_\tau$ match the theoretical bound $M/\tau$.
    \item \textbf{Cooperation from Selfishness}: In multi-NE games,
    increasing search budget improves welfare for \emph{all} agents.
\end{enumerate}


\bibliographystyle{plain}
\begin{thebibliography}{10}

\bibitem{nash1950}
J.~Nash. Equilibrium points in $n$-person games. \emph{PNAS}, 36(1):48--49, 1950.

\bibitem{kakutani1941}
S.~Kakutani. A generalization of Brouwer's fixed point theorem.
\emph{Duke Math. J.}, 8(3):457--459, 1941.

\bibitem{mckelvey1995}
R.~McKelvey and T.~Palfrey. Quantal response equilibria for normal form games.
\emph{Games and Economic Behavior}, 10(1):6--38, 1995.

\bibitem{giannou2022}
A.~Giannou, K.~Lotidis, P.~Mertikopoulos, and P.~Vlatakis-Gkaragkounis.
On the convergence of policy gradient methods to Nash equilibria in general
stochastic games. \emph{NeurIPS}, 2022.

\bibitem{chao1984}
A.~Chao. Nonparametric estimation of the number of classes in a population.
\emph{Scandinavian Journal of Statistics}, 11(4):265--270, 1984.

\bibitem{good1953}
I.~J.~Good. The population frequencies of species and the estimation of
population parameters. \emph{Biometrika}, 40(3/4):237--264, 1953.

\bibitem{foerster2018}
J.~Foerster, R.~Chen, M.~Al-Shedivat, S.~Whiteson, P.~Abbeel, and I.~Mordatch.
Learning with opponent-learning awareness. \emph{AAMAS}, 2018.

\bibitem{kim2021}
D.-K.~Kim, M.~Liu, M.~Riemer, C.~Sun, M.~Abdulhai, G.~Habibi,
S.~Lopez-Cot, G.~Tesauro, and J.~P.~How. A policy gradient algorithm
for learning to learn in multiagent reinforcement learning.
\emph{ICML}, 2021.

\end{thebibliography}

\end{document}
